% Requires amsmath, amssymb, booktabs, graphicx and url (or hyperref).
% Upload the two accompanying vector PDFs into Images/.
\section{Hardware efficiency of 64-sample attention}
\label{sec:fpga_evaluation}

We optimized the physical realization of FROST at a fixed stream length
$L=64$ while preserving its sampled computation. Stage-local state
updates, range-bounded accumulation and simplified completion control
reduced implementation overhead without changing the emitted output
words or execution cycles on the verification workloads. The optimized
kernel used 6,365 combinational resources and 2,036 flip-flops (FFs),
with no digital signal processing (DSP) or block random-access memory
(BRAM) instances on the target field-programmable gate array (FPGA).
In a separate standard-cell assessment, it achieved lower mapped area,
estimated active power and tile latency than each of the three
implemented arithmetic controls.

\paragraph{Preserving the computation while simplifying its implementation.}
The fused kernel receives 16 attention scores and a $16\times16$ value
matrix. Each category emitted by the tensorized finite-state operator
selects one value row for accumulation, eliminating explicit
probability--value multiplication. The operator retains tensor-train
rank 2, 16 finite-state-machine (FSM) states per input, 8-bit cumulative
distribution function (CDF) entries, 16 burn-in updates and 32 separately
initialized xorshift32 random states. Coefficients, random sequences,
score preprocessing and input/output formats are unchanged by the
implementation optimizations.

The first change replaces state-history storage with scheduled local
updates. In the original 16-stage pipeline, the four-bit state for stage
$j\in\{0,\ldots,15\}$ was delayed by $j$ cycles, requiring
$4\sum_{j=0}^{15}j=480$ history bits. The revised pipeline instead delays
each stage's FSM and driving-random-state update enable. A 15-bit enable
delay line replaces the history bank, and each stage consumes its local
state directly. This preserves the association between sample index,
FSM state and tensor-train transition while reducing storage and
associated switching (Fig.~\ref{fig:fused_hardware_stage3}a).

The second change specializes accumulation and control for the fixed
64-sample workload. Let $x_{t,d}$ denote the signed 16-bit integer value
selected at retained sample $t\in\{1,\ldots,64\}$ for output channel
$d\in\{1,\ldots,16\}$. The integer accumulator $S_d$ and rounded output
word $q_d$ satisfy
\begin{equation}
S_d=\sum_{t=1}^{64}x_{t,d},\qquad
-2^{21}\leq S_d\leq 2^{21}-64,\qquad
q_d=\left\lfloor\frac{S_d+32}{64}\right\rfloor.
\label{eq:frost_L64_accumulator}
\end{equation}
The sum, including the rounding offset, fits in 22 signed bits.
Sixteen 22-bit accumulators therefore replace the original 32-bit
accumulators, and the final scaling uses a fixed arithmetic shift.
Output words retain the common signed 32-bit interface with ten
fractional bits. Finally, the trailing edge of the contiguous pipeline
validity sequence identifies the last sample, removing the received-sample
counter; the launch counter is reduced to six bits.

\paragraph{Implementation controls and verification.}
The controls comprise single-precision floating-point (FP32)
implementations with one or four shared fused multiply-add (FMA) lanes,
using Berkeley HardFloat~\cite{hauser2019hardfloat}, and a four-lane
fixed-point implementation (FX16). Both FP32 controls use a 1,025-entry
exponential look-up table (LUT), with spacing $1/64$ over score
differences from 0 to 16. They accumulate unnormalized
exponential--value products and compute one reciprocal per query to
normalize the output vector. FX16 uses a LUT with 16 fractional bits,
widened accumulators and one restoring reciprocal per query.

All implementations accept signed 16-bit scores and values with ten
fractional bits and return signed 32-bit words with ten fractional bits.
FP32 computation operates on the integer-unit value representation
before output conversion. Each tile comprises 16 consecutive queries
sharing one value matrix. Transfers use one input word per cycle;
the 256 value words are loaded once per tile and each query loads
16 scores. Queries do not overlap. Sorting, inverse index mapping,
random-number generation, CDF storage, buffering, burn-in, conversion
and output transfers remain inside the FROST boundary. One-time seed
initialization is excluded.

The optimized design was checked on 3,200 query rows from 200 evaluation
matrices, using stream seeds 4004, 5005 and 6006. All 153,600 output
components agreed bit for bit with both the pre-optimization kernel
and an independent C++ reference; query cycle counts also matched.
An additional 3,072 output comparisons covered full-range signed inputs,
including constant minimum and maximum values, alternating extrema and
random values. Across the three implementation revisions, 460,800
regular-workload output components passed the independent-reference
checks. These tests verify implementation preservation; approximation
error was not matched between FROST and the arithmetic controls.

\paragraph{FPGA resource savings.}
Yosys 0.33 and nextpnr-ecp5 0.6 were used to synthesize and route the
kernels on an LFE5U-85F FPGA with CABGA381 packaging and speed
grade 6~\cite{yosysSoftware,nextpnrSoftware}. All optimized variants
completed routing with seeds 1, 2 and 3 at the common 8-MHz clock.
Resource counts were identical across routing seeds. The unchanged
arithmetic controls were evaluated using their existing implementation
reports from the same device and flow.

Relative to the pre-optimization FROST kernel at $L=64$, the final
implementation reduced combinational resources by 6.9\% and FFs by
24.6\%. Compared with four-lane FP32, it reduced combinational resources
by 73.3\% and eliminated all 16 DSP and three BRAM instances.
Its four reported resource counts were each lower than those of
either FP32 control (Table~\ref{tab:fused_fpga_resources}).
The dominant savings are in combinational logic and dedicated
arithmetic and memory blocks.
FROST stores values in distributed memory and implements constant CDFs
as logic, so the zero-BRAM result includes the cost of alternative
storage resources.

\begin{table}[!htbp]
\centering
\caption{Post-route FPGA resources for the 64-sample attention comparison.}
\label{tab:fused_fpga_resources}
\small
\setlength{\tabcolsep}{4pt}
\begin{tabular*}{\linewidth}{@{\extracolsep{\fill}}lrrrr@{}}
\toprule
Implementation & COMB & FF & DSP & BRAM \\
\midrule
FP32-LUT, 1 lane & 9,304 & 2,043 & 4 & 3 \\
FP32-LUT, 4 lanes & 23,799 & 2,043 & 16 & 3 \\
FX16, 4 lanes & 3,419 & 1,280 & 7 & 2 \\
FROST, before optimization & 6,834 & 2,699 & 0 & 0 \\
FROST, optimized & 6,365 & 2,036 & 0 & 0 \\
\bottomrule
\end{tabular*}
\par\smallskip\footnotesize
Both FROST rows use $L=64$, identical coefficients and identical
input/output formats. COMB denotes ECP5 \texttt{TRELLIS\_COMB}, including
arithmetic and distributed-memory logic. DSP and BRAM denote
\texttt{MULT18X18D} and nominal 18-kbit \texttt{DP16KD} instances.
The counts include all declared kernel support logic. Board-specific
pin placement and external input/output delays were not constrained.
\end{table}

\begin{figure}[!htbp]
\centering
\includegraphics[width=\linewidth]{Images/FROST_hardware_stage3.pdf}
\caption{\textbf{Reducing hardware overhead while preserving 64-sample attention.}
\textbf{a}, Implementation changes: stage-local FSM updates replace
480 bits of state history with a 15-bit update-enable delay line;
range analysis reduces accumulation width from 32 to 22 bits;
pipeline validity replaces received-sample counting. The storage and
width annotations describe logical structures, whereas mapped device
counts are shown in \textbf{b}.
\textbf{b}, Absolute post-route FPGA resources for all arithmetic controls
and FROST before and after optimization. The final FROST implementation
uses fewer COMB and FF resources than both FP32 controls and no DSP or
BRAM instances. Axes are independent across resource classes and begin
at zero. Both FROST implementations use $L=64$.}
\label{fig:fused_hardware_stage3}
\end{figure}

\paragraph{Standard-cell power, performance and area.}
The same RTL variants were mapped to the Nangate45 typical standard-cell
library at 1.10 V and $25\,^{\circ}$C, with memories implemented using
standard cells. OpenSTA 3.1.0 estimated active power from zero-delay
gate-level switching activity at a common 15-MHz
clock~\cite{openstaSoftware}. The activity workload was the same
16-query tile used for the controls, including loading, burn-in,
computation and output transfers. FROST used three stream seeds.
All variants met pre-layout register timing at this frequency.
Across the three optimization revisions, all 2,304 gate-level output
components matched the corresponding RTL simulations, and every power
analysis reported zero unannotated pins.

Let $C_{\mathrm{tile}}$ denote the tile cycle count, $f$ the operating
frequency in hertz, and $P_{\mathrm{cell}}$ the estimated active cell
power in watts, including internal, switching and leakage components.
Tile latency $\tau_{\mathrm{tile}}$ in seconds and estimated energy
$E_{\mathrm{tile}}$ in joules are
\begin{equation}
\tau_{\mathrm{tile}}=\frac{C_{\mathrm{tile}}}{f},\qquad
E_{\mathrm{tile}}=P_{\mathrm{cell}}\tau_{\mathrm{tile}}.
\label{eq:fused_tile_energy}
\end{equation}
All standard-cell comparisons use $f=15$ MHz; FPGA throughput is not
used to calculate these energy values.

Optimization reduced FROST's mapped area from
$63.62\times10^3$ to $58.14\times10^3~\mu\mathrm{m}^2$, an 8.6\%
reduction. Estimated active power decreased from 2.650 to 2.471~mW.
The tile cycle count remained 2,592, corresponding to 172.8~$\mu$s
at 15 MHz or 324~$\mu$s at the common 8-MHz FPGA clock.
Estimated energy per tile decreased from 0.4579 to 0.4271~$\mu$J,
a 6.7\% reduction obtained with unchanged sampled outputs and
execution cycles.

The optimized FROST kernel had the lowest mapped area, estimated active
power and tile latency among the three evaluated arithmetic controls
(Table~\ref{tab:fused_asic_power}; Fig.~\ref{fig:fused_hardware_tradeoffs}).
Relative to four-lane FP32, area, power and latency were reduced by
31.9\%, 30.5\% and 8.0\%, respectively, yielding 36.0\% lower energy.
Relative to FX16, the corresponding reductions were 3.2\%, 13.0\%
and 9.0\%, with 20.8\% lower energy. Relative to single-lane FP32,
estimated tile energy was reduced by 70.7\%.
These simultaneous power, performance and area (PPA) improvements
identify a concrete operating point at which categorical attention
offers an implementation advantage over both floating-point and
fixed-point controls in the same standard-cell flow.

\begin{table}[!htbp]
\centering
\caption{Standard-cell PPA and energy for the fixed 64-sample operating point.}
\label{tab:fused_asic_power}
\fontsize{8}{9.5}\selectfont
\setlength{\tabcolsep}{3pt}
\begin{tabular*}{\linewidth}{@{\extracolsep{\fill}}lrrrr@{}}
\toprule
Implementation
& \shortstack{Area\\($10^3\,\mu\mathrm{m}^2$)}
& \shortstack{Active power\\(mW)}
& \shortstack{Tile latency\\($\mu$s)}
& \shortstack{Energy/tile\\($\mu$J)} \\
\midrule
FP32-LUT, 1 lane & 62.86 & 3.598 & 405.3 & 1.4584 \\
FP32-LUT, 4 lanes & 85.36 & 3.556 & 187.7 & 0.6676 \\
FX16, 4 lanes & 60.08 & 2.839 & 189.9 & 0.5391 \\
FROST, before optimization
& 63.62 & $2.650\pm0.004$ & 172.8 & $0.4579\pm0.0007$ \\
FROST, optimized
& 58.14 & $2.471\pm0.002$ & 172.8 & $0.4271\pm0.0004$ \\
\bottomrule
\end{tabular*}
\par\smallskip\footnotesize
Nangate45, 1.10 V, $25\,^{\circ}$C and 15 MHz for every row.
Both FROST rows use $L=64$. Area is summed mapped-cell area; power
and energy are pre-layout tool estimates. FROST values are means
$\pm$ sample standard deviations across three stream seeds on one
activity tile, not uncertainty across images. Deterministic controls
use that tile once. Clock-tree implementation, extracted interconnect,
time-resolved glitches, pads and external memory are excluded.
\end{table}

\begin{figure}[!htbp]
\centering
\includegraphics[width=\linewidth]{Images/FROST_hardware_tradeoffs.pdf}
\caption{\textbf{Simultaneous PPA gains at the fixed 64-sample operating point.}
\textbf{a}, Percentage reductions for optimized FROST relative to each
arithmetic control within the same Nangate45, 15-MHz flow.
Each entry is calculated as $100(1-m_{\mathrm{FROST}}/m_{\mathrm{control}})$,
where $m$ denotes the indicated area, power, latency or energy metric.
Positive values indicate a reduction; shading uses the common
percentage scale. Calculations use unrounded results.
\textbf{b}, Estimated energy per tile before and after the
output-preserving implementation changes. Both FROST versions retain
64 samples and the same tile cycle count. Error bars show sample
standard deviations over three stream seeds on one activity tile and
can be smaller than the plotted marks. A tile comprises 16 queries;
all power and energy values are pre-layout estimates.}
\label{fig:fused_hardware_tradeoffs}
\end{figure}

The evaluation covers the fused score-to-attention-output kernel,
including its numerical support operations. Query--key products,
projections, feed-forward layers and whole-network classification
remain outside this boundary. The results establish implementation-level
resource and energy savings for the stated operating point; the power
values are standard-cell estimates rather than FPGA board measurements.
